% ══════════════════════════════════════════════════════════════════════════════
%  CHAPTER 2: FRIEDMANN EQUATIONS & COSMOLOGICAL DYNAMICS
% ══════════════════════════════════════════════════════════════════════════════

\chapter{Friedmann Equations \& Cosmological Dynamics}

% ──────────────────────────────────────────────────────────────────────────────
\section{The FLRW Metric}

\subsection{Physics Intuition}

\begin{intuition}
The principle of \emph{cosmic homogeneity and isotropy}—that on large scales the
universe looks the same from every point and in every direction—drastically
simplifies the mathematical description of spacetime.
This is not merely a convenient assumption: CMB temperature maps show variations
of only one part in $10^5$, and galaxy surveys above $\sim100\,\mathrm{Mpc}$ are
remarkably uniform. These two symmetries together constrain the spacetime metric
to the unique Friedmann--Lema\^itre--Robertson--Walker (FLRW) form.
\end{intuition}

\subsection{Theory}

\begin{theory}
The \textbf{FLRW metric} describes a homogeneous, isotropic universe:
\begin{equation}
  ds^2 = -c^2\,dt^2 + a(t)^2
    \!\left[\frac{dr^2}{1-kr^2} + r^2\,d\theta^2 + r^2\!\sin^2\!\theta\,d\phi^2\right],
\end{equation}
with curvature parameter $k\in\{-1,0,+1\}$.

The curvature parameter $k$ determines the global geometry of spatial slices:
$k=+1$ for a closed (sphere-like) universe, $k=-1$ for an open (saddle-like)
universe, and $k=0$ for a spatially flat (Euclidean) universe.
The comoving coordinate $r$ labels a fixed point in the matter distribution;
physical (proper) distances scale with the scale factor as $d_{\rm phys} = a(t)\,r$.
Current Planck observations constrain $|\Omega_k| < 0.005$, consistent with exact flatness.
\end{theory}

\subsection{Question Examples}

\begin{example}{Proper Distance in FLRW}{}
Two galaxies are separated by a comoving distance $r = 1000\,\mathrm{Mpc}$
in a flat ($k=0$) FLRW universe.
Find their proper separation (a)~today and (b)~at $z=2$, when $a(t_e)=1/3$.

\medskip\textbf{Solution.}\quad
(a)~Today $a(t_0)=1$, so $d_{\rm phys} = 1000\,\mathrm{Mpc}$.
(b)~At $z=2$, $d_{\rm phys} = a(t_e)\cdot r = (1/3)\times1000 \approx 333\,\mathrm{Mpc}$;
the universe was three times smaller and the galaxies three times closer together.
\end{example}

\subsection{Extra}

\begin{extra}
The FLRW metric was independently derived by Alexander Friedmann (1922),
Georges Lema\^itre (1927), Howard Robertson (1935), and Arthur Walker (1936).
Lema\^itre was the first to combine it with recession-velocity data to argue
that the universe is expanding—two years before Hubble's 1929 paper—and
also extrapolated backwards to propose what he called the \emph{primeval atom},
now known as the Big Bang.
\end{extra}

% ──────────────────────────────────────────────────────────────────────────────
\section{The Friedmann Equations}

\subsection{Physics Intuition}

\begin{intuition}
The Friedmann equations are the equations of motion for the universe as a whole.
The first equation resembles Newtonian energy conservation: the square of the expansion
rate ($H = \dot a/a$) equals the total energy density divided by an appropriate factor.
The second (acceleration) equation tells us whether expansion speeds up or slows down:
ordinary matter and radiation have positive pressure and \emph{decelerate} expansion,
while a cosmological constant (dark energy) with $w=-1$ \emph{accelerates} it.
\end{intuition}

\subsection{Theory}

\begin{theory}
Inserting the FLRW metric into the Einstein field equations gives the
\textbf{Friedmann equations}:
\begin{align}
  \left(\frac{\dot{a}}{a}\right)^{\!2}
    &= \frac{8\pi G}{3}\rho - \frac{kc^2}{a^2} + \frac{\Lambda c^2}{3},
    \label{eq:friedmann1}\\[6pt]
  \frac{\ddot{a}}{a}
    &= -\frac{4\pi G}{3}\!\left(\rho + \frac{3p}{c^2}\right) + \frac{\Lambda c^2}{3}.
    \label{eq:friedmann2}
\end{align}
The equation-of-state parameter $w = p/(\rho c^2)$ characterises each component:
$w=0$ for pressureless dust (cold dark matter or baryons), $w=1/3$ for radiation
and relativistic species, and $w=-1$ for a cosmological constant.
Components with $w < -1/3$ contribute positively to the acceleration
Eq.~\eqref{eq:friedmann2}, driving accelerated expansion.
Combining both Friedmann equations yields the \emph{fluid equation},
$\dot\rho + 3H(\rho + p/c^2) = 0$, expressing covariant energy conservation for
each component independently.
\end{theory}

\subsection{Question Examples}

\begin{example}{Scale Factor of a Flat Matter-Dominated Universe}{}
Derive the time dependence $a(t)$ for a flat, radiation-dominated universe
($k=0$, $\Lambda=0$, $w=1/3$).

\medskip\textbf{Solution.}\quad
For $k=0$, $\Lambda=0$ and pressureless dust ($p=0$), Eq.~\eqref{eq:friedmann1}
reduces to $\dot{a}^2 \propto a^{-1}$, giving:
\begin{equation}
  a(t) \;\propto\; t^{2/3}.
\end{equation}
For radiation ($w=1/3$), $\rho\propto a^{-4}$, so $\dot a^2 \propto a^{-2}$ and
$a \propto t^{1/2}$—a slower expansion than the matter-dominated case.
\end{example}

\subsection{Extra}

\begin{extra}
Friedmann published his expanding-universe solution in 1922; Einstein initially
published a rebuttal claiming a mathematical error, then retracted it.
It was not until Hubble's 1929 observational evidence that the scientific community
accepted the expanding universe.
The steady-state theory—championed by Hoyle, Gold, and Bondi until the
1960s—proposed continuous matter creation to maintain constant mean density,
but was decisively ruled out by the discovery of the cosmic microwave background
in 1965.
\end{extra}

\clearpage
